% ===============================================================
% SL-06d-ML: Stundenleistung MUSTERLÖSUNG
% Spanisch Klasse 7 - Sequenz 6: Mi ciudad
% ===============================================================

\documentclass[11pt, a4paper]{article}

\newif\ifswmodus
\swmodusfalse

\usepackage[utf8]{inputenc}
\usepackage[T1]{fontenc}
\usepackage[ngerman]{babel}
\usepackage{helvet}
\renewcommand{\familydefault}{\sfdefault}

\usepackage[margin=2cm]{geometry}
\usepackage{enumitem}
\usepackage{tabularx}
\usepackage{booktabs}

\usepackage{xcolor}
\usepackage{tcolorbox}
\tcbuselibrary{skins}

\usepackage{fancyhdr}
\usepackage{lastpage}
\usepackage{needspace}

\definecolor{spanisch-color}{HTML}{c41e3a}
\definecolor{grau-color}{HTML}{888888}
\definecolor{hellgrau-color}{HTML}{f5f5f5}
\definecolor{loesung-color}{HTML}{c62828}

\definecolor{sw-dark}{HTML}{000000}
\definecolor{sw-gray}{HTML}{666666}
\definecolor{sw-white}{HTML}{ffffff}

\ifswmodus
    \colorlet{spanisch}{sw-dark}
    \colorlet{grau}{sw-gray}
    \colorlet{hellgrau}{sw-white}
    \colorlet{loesung}{sw-dark}
\else
    \colorlet{spanisch}{spanisch-color}
    \colorlet{grau}{grau-color}
    \colorlet{hellgrau}{hellgrau-color}
    \colorlet{loesung}{loesung-color}
\fi

\newcommand{\fachfarbe}{spanisch}

\newcommand{\thema}{Stundenleistung: Mi ciudad}
\newcommand{\fach}{Spanisch}
\newcommand{\klasse}{7}
\newcommand{\maxpunkte}{20}
\newcommand{\zeit}{25 Minuten}
\newcommand{\datum}{\today}

\pagestyle{fancy}
\fancyhf{}
\renewcommand{\headrulewidth}{0.4pt}
\renewcommand{\footrulewidth}{0.4pt}

\lhead{\fach{} -- Klasse \klasse}
\chead{\textcolor{loesung}{\textbf{SL-06d MUSTERLÖSUNG}}}
\rhead{\datum}

\lfoot{\zeit}
\cfoot{}
\rfoot{Seite \thepage\ von \pageref{LastPage}}

\newcommand{\punktebox}[1]{\hfill\fbox{\small \_/#1}}
\newcommand{\luecke}[1]{\rule{#1}{0.4pt}}
\newcommand{\lsg}[1]{\textcolor{loesung}{\textbf{#1}}}

\newtcolorbox{infobox}{
    colback=hellgrau,
    colframe=\fachfarbe,
    boxrule=1pt,
    arc=0pt,
    left=8pt, right=8pt, top=8pt, bottom=8pt
}

\newenvironment{aufgabe}[1]{%
    \needspace{5\baselineskip}%
    \nopagebreak[4]%
    \vspace{10pt}
    \noindent\textbf{\textcolor{\fachfarbe}{Aufgabe #1}}
    \nopagebreak[4]%
    \vspace{5pt}
    \nopagebreak[4]%
    \begin{list}{}{\leftmargin=0pt}
    \item[]
}{%
    \end{list}
}

\newcommand{\es}[1]{\textcolor{\fachfarbe}{\textbf{#1}}}

\begin{document}

\begin{center}
    {\Large\bfseries\textcolor{\fachfarbe}{\thema}}\\[0.3em]
    {\large\textcolor{loesung}{\textbf{--- MUSTERLÖSUNG ---}}}
\end{center}

\vspace{5pt}

\noindent
\begin{tabularx}{\textwidth}{|X|X|X|}
\hline
\textbf{Name:} \lsg{Musterschüler/in} &
\textbf{Klasse:} \klasse &
\textbf{Datum:} \datum \\
\hline
\end{tabularx}

\vspace{10pt}

\begin{infobox}
\textbf{Zeit:} \zeit{} | \textbf{Punkte:} \maxpunkte{} | \textbf{Hilfsmittel:} keine
\end{infobox}

\vspace{15pt}

% ===============================================================
% AUFGABE 1: Zuordnung Orte
% ===============================================================

\begin{aufgabe}{1 -- Orte in der Stadt zuordnen}
Verbinde die spanischen Orte mit der deutschen Bedeutung. \punktebox{4}
\end{aufgabe}

\begin{center}
\begin{tabular}{rcl}
\es{el supermercado} & \lsg{---} & das Krankenhaus \lsg{$\rightarrow$ supermercado = Supermarkt} \\[0.8em]
\es{la estación} & \lsg{---} & der Supermarkt \lsg{$\rightarrow$ estación = Bahnhof} \\[0.8em]
\es{el hospital} & \lsg{---} & der Platz \lsg{$\rightarrow$ hospital = Krankenhaus} \\[0.8em]
\es{la plaza} & \lsg{---} & der Bahnhof \lsg{$\rightarrow$ plaza = Platz} \\
\end{tabular}
\end{center}

\vspace{15pt}

% ===============================================================
% AUFGABE 2: ir und estar Konjugation
% ===============================================================

\begin{aufgabe}{2 -- Die Verben \es{ir} und \es{estar}}
Ergänze die richtige Form des Verbs. \punktebox{4}
\end{aufgabe}

\noindent
a) Yo \lsg{voy} (ir) al cine. \hfill (ich gehe)

\vspace{0.8em}

\noindent
b) Tú \lsg{estás} (estar) en la plaza. \hfill (du bist)

\vspace{0.8em}

\noindent
c) Ella \lsg{va} (ir) a la tienda. \hfill (sie geht)

\vspace{0.8em}

\noindent
d) Nosotros \lsg{estamos} (estar) en el parque. \hfill (wir sind)

\vspace{15pt}

% ===============================================================
% AUFGABE 3: Wegbeschreibung Lückentext
% ===============================================================

\begin{aufgabe}{3 -- Wegbeschreibung verstehen}
Wähle das richtige Wort aus der Box und ergänze. \punktebox{6}
\end{aufgabe}

\begin{center}
\fbox{\parbox{0.8\textwidth}{\centering
\es{todo recto} \quad \es{a la derecha} \quad \es{a la izquierda}\\[0.3em]
\es{cerca de} \quad \es{enfrente de} \quad \es{al lado de}
}}
\end{center}

\vspace{0.8em}

\noindent
a) -- Perdona, ¿dónde está el museo?\\
\hspace*{0.5cm} -- Sigue \lsg{todo recto} y luego gira \lsg{a la izquierda}. (geradeaus, nach links)

\vspace{0.8em}

\noindent
b) -- ¿Está lejos el banco?\\
\hspace*{0.5cm} -- No, está \lsg{al lado de} la iglesia. (neben)

\vspace{0.8em}

\noindent
c) -- ¿Dónde está la farmacia?\\
\hspace*{0.5cm} -- Está \lsg{enfrente de} el hospital. (gegenüber von)

\vspace{15pt}

% ===============================================================
% AUFGABE 4: Wegbeschreibung schreiben
% ===============================================================

\begin{aufgabe}{4 -- Einen Weg beschreiben}
Du bist am Bahnhof (\es{la estación}). Jemand fragt nach dem Weg zum Kino (\es{el cine}).\\
Schreibe einen kurzen Dialog mit 3 Sätzen auf Spanisch. \punktebox{6}
\end{aufgabe}

\noindent
\textbf{Turista:} Perdona, ¿dónde está el cine?

\vspace{0.8em}

\noindent
\textbf{Tú:} \lsg{Sigue todo recto y gira a la derecha. / El cine está cerca de aquí.} \hfill (2P)

\vspace{1.2em}

\noindent
\textbf{Turista:} ¿Está lejos?

\vspace{0.8em}

\noindent
\textbf{Tú:} \lsg{No, está cerca. / No, está al lado del parque. / Está a cinco minutos.} \hfill (2P)

\vspace{1.2em}

\noindent
\textbf{Turista:} ¡Muchas gracias!

\vspace{0.8em}

\noindent
\textbf{Tú:} \lsg{¡De nada! / ¡Adiós!} \hfill (2P)

% ===============================================================
% PUNKTEÜBERSICHT
% ===============================================================

\vspace{25pt}
\hrule
\vspace{10pt}

\noindent
\begin{tabularx}{\textwidth}{|X|c|c|c|c||c|}
\hline
\textbf{Aufgabe} & 1 & 2 & 3 & 4 & \textbf{Gesamt} \\
\hline
\textbf{max. Punkte} & 4 & 4 & 6 & 6 & \textbf{20} \\
\hline
\textbf{erreicht} & \lsg{4} & \lsg{4} & \lsg{6} & \lsg{6} & \lsg{20} \\
\hline
\end{tabularx}

\vspace{10pt}

\begin{flushright}
\textbf{Note:} \lsg{1}
\end{flushright}

\vspace{1em}

\textbf{Notenschlüssel:}

\begin{tabular}{|c|c|c|c|c|c|c|}
\hline
\textbf{Note} & 1 & 2 & 3 & 4 & 5 & 6 \\
\hline
\textbf{ab Punkte} & 19 & 16 & 12 & 8 & 4 & <4 \\
\hline
\end{tabular}

\vspace{1em}

\textbf{Bewertungshinweise:}
\begin{itemize}
    \item Aufgabe 1: Je 1P pro richtige Zuordnung
    \item Aufgabe 2: Je 1P pro richtige Verbform
    \item Aufgabe 3: Je 1P pro richtige Lücke (6 Lücken)
    \item Aufgabe 4: Inhalt (3P) + Sprache (3P)
    \begin{itemize}
        \item Wegbeschreibung verständlich: 2P
        \item Antwort auf ``lejos'': 2P
        \item Verabschiedung: 2P
        \item Bei Rechtschreibfehlern: -0,5P pro Fehler (max. -2P gesamt)
    \end{itemize}
\end{itemize}

\end{document}
